\section{Synthesising States and Deriving the Reducer}
\label{sec:derivation}

The source reduction results provide the starting point for our maintained representation.
Expanding their definitions into context-independent reductions gives a reference state tuple whose components already have joins.
We use this tuple to specify the information that auxiliary states must supply: together with the retained source states, they must recover every reference component.
Synthesis therefore searches only for auxiliary states.
Once a recovery map is constructed and verified, the join follows by composition.
This section develops that process through a Householder-type norm, where the implementation uses synthesis to select an auxiliary state.

\subsection{Source States and Reference Representation}
\label{sec:selection}

Let $\stsrc(S)$ denote the source-derived initial state tuple on a legal non-empty part $S$.
These states retain the source definitions whenever they are defined on legal parts.
A demonstrated domain defect is repaired by extracting the offending context: its coefficients replace that reduction in $\stsrc$, and the source result is reconstructed for the outputs.
We then expand these definitions using the context-separation rules, substitute the definitions of their contexts, and collect the distinct context-independent reductions.
Their values form the reference tuple $\stref(S)$, and the expanded expressions give a known map $\Gsrc$:
\begin{equation}
\label{eq:source-reference}
\stsrc(S)=\Gsrc(\stref(S)).
\end{equation}
The reference components join independently and remain fixed throughout auxiliary-state search.
Whereas the existence analysis starts from the required outputs, this construction starts from the source states that the fused program will retain.

For the running example, consider two input sequences $x_j,e_j$.
A Householder-type norm first computes a squared norm $q_S$ and then uses $k_S=\sqrt{q_S}$ as the context of a second reduction:
\[
q_S=\sum_{j\in S}x_j^2,\qquad
v_S=\sum_{j\in S}(x_j+k_Se_j)^2.
\]
Thus $\stsrc(S)=(q_S,v_S)$; the scalar computation $k_S$ adds no reduction state.
Expanding the second definition exposes two further context-independent quantities,
\[
u_S=\sum_{j\in S}e_j^2,\qquad
p_S=\sum_{j\in S}x_je_j,\qquad
v_S=q_S+2k_Sp_S+q_Su_S.
\]
We obtain $\stref(S)=(q_S,u_S,p_S)$ and
$\Gsrc(q,u,p)=(q,q+2\sqrt q\,p+qu)$.
The two source states leave the information in $u$ and $p$ combined; auxiliary-state synthesis determines what additional information to retain.

\subsection{Discovering and Synthesising Auxiliary States}
\label{sec:synthesis}

Candidate expressions come from how a retained result changes with its context.
The analyser symbolically shifts or scales the local contexts and separates the resulting terms.
For a sum state $r$, a successful decomposition under the algebraic rules has the form
\begin{equation}
\label{eq:features}
r(S;\kappa')=\sum_{\ell}\alpha_\ell(\delta,\rho)\,e_\ell(S;\kappa_S),
\end{equation}
where $\kappa_S$ is the local context, $\kappa'$ the target context, and $\delta,\rho$ describe their additive or multiplicative changes.
The factors $\alpha_\ell$ depend on these changes; the aggregate coefficients $e_\ell$ expose information needed to correct the retained result.
Terms with different correction factors can also provide separate candidates.
For extrema, the corresponding monotonicity and context-extraction rules serve this role.
Each candidate has a definition in $\stref$ for reasoning and a structural reduction expression retaining local contexts for code generation.
The retained candidates form $\mathcal T_{\mathrm{pre}}$; each is a defined function of the reference tuple and fixed program parameters.

In the Householder example, writing $v_S(k)=\sum_{j\in S}(x_j+ke_j)^2$ gives
\[
v_S(k_S+\delta)
 =v_S+2\delta(p_S+k_Su_S)+\delta^2u_S.
\]
Besides the retained $v_S$, this correction exposes $p_S+k_Su_S$ and $u_S$.
Adding both gives four maintained states, while the existence construction supplies a default limit of three.
The search must therefore select a smaller augmentation.

Let $C=(c_1,\ldots,c_w)$ be the auxiliary expressions to be synthesised.
For a reference value $\xi$, their definitions determine
\[
G(\xi)=(\Gsrc(\xi),C(\xi)),\qquad
\st(S)=G(\stref(S)).
\]
We require $\st$ to recover $\stref$, a sufficient condition for constructing a join for the entire maintained tuple.
The synthesis specification expresses this information requirement before recovery is constructed.
Let $\mathcal R$ conservatively contain all reachable reference tuples, constrained by finite-sum facts and proved definedness conditions for $\Gsrc$ and the candidates.
For fixed program parameters, synthesis requires
\begin{equation}
\label{eq:synthesis}
\forall \xi,\eta\in\mathcal R:\quad
\bigl(\Gsrc(\xi)=\Gsrc(\eta)\ \land\ C(\xi)=C(\eta)\bigr)
\ \Longrightarrow\ \xi=\eta.
\end{equation}
Auxiliary states must distinguish reference tuples that the initial states cannot distinguish.
The known reference join then supplies the basis for joining the selected representation, so no join expression is included among the synthesis unknowns.

The candidate expressions in a set $\mathcal T$ serve as grammar terminals. Each auxiliary expression is drawn from the finite grammar
\[
\begin{aligned}
\mathcal E_1&::=e\quad(e\in\mathcal T),\\
\mathcal E_\ell&::=\mathcal E_{\ell-1}
 \mid\mathcal E_{\ell-1}+\mathcal E_1
 \mid\mathcal E_{\ell-1}-\mathcal E_1
 \qquad(2\le\ell\le d).
\end{aligned}
\]
\begingroup
\setlength{\intextsep}{4pt}
\reservewrapspace{21\baselineskip}
\begin{wrapfigure}{r}{0.48\textwidth}
\begin{minipage}{\linewidth}
\begin{algorithm}[H]
\small
\SetInd{0.5em}{0.7em}
\caption{\textsc{Select}$(\mathit{chain},b_{\max})$}
\label{alg:select}
$\stsrc\leftarrow$ source states after domain repair\;
$(\stref,\Gsrc)\leftarrow$ expand definitions of $\stsrc$\;
\lIf{preparation fails}{\Return keep source chain}
$\mathcal T_{\mathrm{pre}},K\leftarrow$ correction candidates\;
\ForEach{$M\in[\stsrc,(\stsrc,K)]$ fitting $b_{\max}$}{
\lIf{\textsc{Verify}$(M,\stref)$ succeeds}{\Return reducer}
}
\ForEach{$\mathcal T\in[\mathcal T_{\mathrm{pre}},\mathcal T_{\mathrm{mix}}]$}{
\ForEach{admissible auxiliary width $w$}{
\While{proposals and resources remain}{
Propose $C\in\mathcal E_d^{\,w}$\;
\lIf{\textsc{Verify}$((\stsrc,C),\stref)$ succeeds}{\Return reducer}
}
}}
\Return keep source chain\;
\end{algorithm}
\end{minipage}
\end{wrapfigure}
With the default $d=3$, expressions combine at most three occurrences of candidate terms.
A term can already be nonlinear in the reference values, so the grammar combines quantities discovered by analysis while preserving their structural forms.
The first grammar uses $\mathcal T_{\mathrm{pre}}$; the second also includes the reference components as candidate expressions.
When the existence analysis returns $\mathcal Q$, its bound $|\mathcal Q|$ supplies the default maintained-state limit $b_{\max}$.
The supplementary material establishes a general sufficient augmentation; the search here seeks a compact representation within this limit.

\Cref{alg:select} first tests $\stsrc$ alone, then its augmentation with the tuple $K$ of admissible correction coefficients.
It invokes synthesis if these choices do not yield a verified reducer within the state limit.
Here $\mathcal T_{\mathrm{mix}}$ contains $\mathcal T_{\mathrm{pre}}$ and the reference components; the auxiliary width satisfies $1\le w\le b_{\max}-|\stsrc|$.
For each proposal, \textsc{Verify} constructs recovery expressions, proves their identities and domains, and derives and verifies the complete reducer.
Search scheduling and solver limits are described with the implementation.

In the recorded Householder derivation, the initial tuple does not suffice and the complete coefficient augmentation exceeds the three-state limit.
SyGuS selects $C(q,u,p)=u$, giving $\st=(q,v,u)$.
The alternative coefficient $p+ku$ loses $u$ when $q=0$: then $p=k=v=0$ while $u$ can vary.
A later join can make the norm context non-zero, so these parts require different corrections.
The information requirement distinguishes these choices.
Selecting $u$ leaves only $p$ to be recovered, including at zero norm.
\finishwrapspace
\endgroup

\subsection{Constructing and Verifying Recovery}
\label{sec:atoms}
\label{sec:recovery}

For a proposed maintained tuple, the selected definitions give the known map $G$.
Recovery constructs an executable map $H$ from the equations defining that tuple, with
\begin{equation}
\label{eq:recovery}
H(G(\stref(S)))=\stref(S).
\end{equation}
The procedure substitutes known expressions and isolates unresolved reference components using supported algebraic rules.
Repeated substitution expresses every component in maintained states and fixed parameters.
Symbolic normalisation and SymPy establish identities, while Z3 proves domains and guarded implications.
This construction of $H$ is separate from the synthesis of $C$.

For Householder, $q$ and $u$ are retained directly.
Solving the definition of $v$ gives the remaining component and hence recovery:
\[
p=
\begin{cases}
0,&q=0,\\[2pt]
\dfrac{v-q-qu}{2\sqrt q},&q>0,
\end{cases}
\qquad H(q,v,u)=(q,u,p).
\]
\reservewrapspace{16\baselineskip}
% Both panels and their common caption form one full-width figure.
% Shared unchanged by the English and Chinese prose editions.
\par\medskip
\noindent\begin{minipage}{\linewidth}
\captionsetup{type=figure}
\centering
\begin{minipage}[t]{0.48\linewidth}
\vspace{0pt}
\centering
{\small\sffamily\bfseries (a) State relations}\par\medskip
% Shared vector source for the English and Chinese prose editions.
\begin{tikzpicture}[
  summary/.style={draw=rulegray,rounded corners=2pt,fill=boxgray!45,
    text width=2.45cm,minimum height=1.05cm,inner sep=4pt,
    align=center,font=\small},
  arrow/.style={-{Stealth[length=2mm]},semithick},
  label/.style={font=\small,fill=white,inner sep=2pt,align=center}]
\node[summary,draw=origblue] (original) at (0,2.2)
  {$A_0(S),\ A_0(T)$\\[3pt]$(t_p,ss_p)$};
\node[summary,draw=auxpurple] (maintained) at (0,0)
  {$A(S),\ A(T)$\\[3pt]$(t_p,ss_p,\aux{n_p})$};
\node[summary,draw=fixgreen] (merged) at (3.65,0)
  {$A(U)$\\[3pt]$(t_U,ss_U,\aux{n_U})$};
\node[summary] (reference) at (0,-2.6)
  {$B(S),\ B(T)$\\[3pt]$(t_p,q_p,n_p)$};
\node[summary] (combined) at (3.65,-2.6)
  {$B(U)$\\[3pt]$(t_U,q_U,n_U)$};
\draw[arrow,draw=auxpurple] (original) --
  node[label,right,text=auxpurple] {add $n_p=\sum_{j\in p}1$\\\S\ref{sec:selection}--\ref{sec:synthesis}} (maintained);
\draw[arrow,draw=fixgreen] (maintained) --
  node[label,above,text=fixgreen] {$\Mrg$}
  node[label,below=6mm,text=fixgreen] {\S\ref{sec:correctness}--\ref{sec:rebase}} (merged);
\draw[arrow] (maintained) -- node[label,left] {$H$\\\S\ref{sec:recovery}} (reference);
\draw[arrow] (reference) -- node[label,above] {$\oplus$}
  node[label,below] {\S\ref{sec:correctness}} (combined);
\draw[arrow] (combined) -- node[label,right] {$G$\\\S\ref{sec:atoms}} (merged);
\end{tikzpicture}

\end{minipage}\hfill
\begin{minipage}[t]{0.48\linewidth}
\vspace{0pt}
\small
\centerline{\sffamily\bfseries (b) Derivation}
\medskip
On each part, let $t=\sum_jx_j$, $q=\sum_jx_j^2$ and $n=\sum_j1$.
Expanding the source definition gives
\[
G(t,q,n)=(t,q-t^2/n,n).
\]
From $(t,ss)$, recovery leaves $q=ss+t^2/n$ dependent on the unknown $n$.
Adding $n$ gives the complete recovery
\[
H(t,ss,n)=(t,ss+t^2/n,n).
\]
Since $n>0$, both maps are defined and $H\circ G=\mathrm{id}$.
The reference join $\oplus$ adds coordinates.
The upper arrow follows by composition:
\[
\Mrg(a,b)=G\bigl(H(a)\oplus H(b)\bigr).
\]
With $t_U=t_S+t_T$ and $n_U=n_S+n_T$, rewriting yields the centred form
\[
\begin{aligned}
ss_U={}&ss_S+ss_T\\
&+\frac{n_Sn_T}{n_U}
  \left(\frac{t_S}{n_S}-\frac{t_T}{n_T}\right)^2.
\end{aligned}
\]
The emitted join evaluates this correction without materialising $q$.
\end{minipage}
\caption{Variance: auxiliary states, recovery and join derivation.
Here $p\in\{S,T\}$ and $U=S\dunion T$ for disjoint non-empty parts.
Arrow annotations identify the subsections that develop each step:
feature construction and synthesis (\S\ref{sec:selection}--\ref{sec:synthesis}), information equivalence (\S\ref{sec:recovery}),
join derivation (\S\ref{sec:correctness}), and numerical formulation (\S\ref{sec:rebase}).}
\label{fig:recovery-square}
\end{minipage}\par\medskip

The zero branch follows from $q=\sum_jx_j^2=0$: each $x_j$ is zero, so $p=\sum_jx_je_j=0$.
This is an instance of a finite-sum rule used by the verifier: a zero sum of non-negative terms forces each term to vanish.
Both branches recover the reference value on their domains, and a lazy conditional evaluates only the selected branch.
The resulting identity $H\circ G=\mathrm{id}$ establishes complete recovery, including legal zero-norm parts.

\Cref{fig:recovery-square} illustrates the construction: recover each part's $(q,u,p)$, add components, and apply $G$ to obtain the maintained tuple of the union.
Initialisation evaluates the state definitions on a singleton.
\finishwrapspace

\subsection{Deriving the Join and Proving Correctness}
\label{sec:correctness}

The reference components combine independently under $\oplus$, giving
$\stref(S\dunion T)=\stref(S)\oplus\stref(T)$ for disjoint non-empty parts.
Recovery therefore yields the maintained join by composition:
\begin{equation}
\label{eq:merge}
\Mrg(a,b)=G\bigl(H(a)\oplus H(b)\bigr).
\end{equation}

The output function $\epi$ follows the source program's scalar dependencies expressed in selected states; a retained output such as $v$ reads its own component.
Verification checks this reconstruction's identity and domain.

\begin{theorem}[Correctness over $\R$]
\label{thm:correct}
Let $\stref$ be defined on legal non-empty parts and join componentwise, and let $\st=G\circ\stref$.
Suppose $G$ and $H$ are defined on their reachable arguments and $H\circ G$ is the identity on reachable reference tuples.
If $\epi$ reconstructs the source outputs from $\st(\Ax)$ on their source domain, the derived reducer computes these outputs for any partition into non-empty parts and any binary join tree over them.
\end{theorem}
\begin{proof}
For disjoint non-empty parts,
\[
\begin{aligned}
\Mrg(\st(S),\st(T))
 &=G\bigl(H(G(\stref(S)))\oplus H(G(\stref(T)))\bigr)\\
 &=G\bigl(\stref(S)\oplus\stref(T)\bigr)\\
 &=G(\stref(S\dunion T))=\st(S\dunion T).
\end{aligned}
\]
Recovery gives the second equality and componentwise aggregation the third.
Initialisation supplies the singleton states, so induction over the tree yields $\st(\Ax)$ at the root.
Applying the verified output reconstruction completes the proof.
\end{proof}

This invariant covers every maintained component, including the auxiliary states.
The associativity and commutativity established in the reducer contract therefore apply to the derived join.
Complete reference recovery provides one common condition for selecting states and constructing their join.

\subsection{Numerical Formulation of the Join}
\label{sec:rebase}

The compositional join specifies the result to compute.
Its direct evaluation can reconstruct expanded reference quantities before recomputing maintained values; such intermediate subproblems can amplify error even when the overall computation is well conditioned~\citep{higham_decomposition}.
For execution, the context-change decompositions used in candidate discovery instead give formulas that correct the retained partial results to the joined context.
Recovery expresses the required coefficients in maintained states, and the corrected results combine under the source reduction operator.

For Householder, let $U=S\dunion T$, $k_P=\sqrt{q_P}$ and $\delta_P=k_U-k_P$ for $P\in\{S,T\}$.
The join adds $q$ and $u$ componentwise and updates the retained $v$ by
\[
v_U=\sum_{P\in\{S,T\}}
 \left[v_P+2\delta_P(p_P+k_Pu_P)+\delta_P^2u_P\right],
\]
where $p_P$ is recovered by the guarded expression above.
The formula retains each part's computed $v_P$ and adjusts it for the changed context.
Centred-moment joins and online softmax similarly correct partial results by recentring or rescaling~\citep{pebay,milakov}.

The verifier proves equality with \eqref{eq:merge} and definedness of the executed branches.
Components without a supported correction use their compositional expressions.
\begin{theorem}[Correctness of the emitted join]
\label{thm:printed}
An emitted join verified equal to the compositional join on reachable states, with its executed branches defined, satisfies the same join invariant.
\end{theorem}
\begin{proof}
Substitution of the verified equality preserves the invariant proved above.
\end{proof}
The emitted program then undergoes the static and differential numerical checks that determine floating-point acceptance.
